\section{Direkte Faltung in der Zeitdomäne}

Der folgende Quellcodeausschnitt zeigt eine Implementierung der direkten, diskreten Faltung
nach Summenformel:
\begin{equation}
y[n] = (x * h)[n] = \sum_{k=0}^{M-1} h[k] \cdot x[n-k]
\label{eq:directConvolution}
\end{equation}

% Farben für Code definieren
\definecolor{codegreen}{rgb}{0,0.6,0}
\definecolor{codegray}{rgb}{0.5,0.5,0.5}
\definecolor{codepurple}{rgb}{0.58,0,0.82}
\definecolor{backcolour}{rgb}{0.95,0.95,0.92}

\lstset{
    language=C++,
    backgroundcolor=\color{backcolour},   
    commentstyle=\color{codegreen},
    keywordstyle=\color{blue},
    numberstyle=\tiny\color{codegray},
    stringstyle=\color{codepurple},
    basicstyle=\ttfamily\small,
    breakatwhitespace=false,         
    breaklines=true,                 
    captionpos=b,                    
    keepspaces=true,                 
    numbers=left,                    
    numbersep=5pt,                  
    showspaces=false,                
    showstringspaces=false,
    showtabs=false,                  
    tabsize=2
}

\lstinputlisting
[language=C++, linerange={20-38}, caption={Grundalgorithmus der Faltung im Zeitbereich},
label={lst:directConvolution}
]
{code/test/JuceTest/Source/convolution.cpp}

Als Parameter des Algorithmus \ref{lst:directConvolution} werden zwei Listen, nämlich jeweils die Folgen
an normalisierten Abttaswerten des Eingangssignals und der Impulssantwort übergeben.
Dadurch, dass die Beträge der Folgen bekannt sind, lässt sich in Zeilen 5 bis einschl. 8
im Voraus die Länge der Ausgangssignalfolge nach $|y[n]| = |h[n]| + |x[n]| - 1$ berechnen und die Werte mit 0 initialisieren.
Schließlich wird die Formel \ref{eq:directConvolution} in Form der verschachtelten Schleife
in Zeilen 9 bis 16 umgesetzt. In manchen Schleifeniterationen kommt es dazu, dass der Index von $x[n - k]$
kleiner als null wird oder die Länge des Eingangssignals übersteigt. Diese Fälle werden in Zeile 11 geprüft und
ggf. übersprungen.